% Vietnamese translation for OI Public Library
% Source-PDF: IOI中国国家候选队论文/国家集训队2018论文集.pdf
\documentclass[11pt]{article}
\usepackage{oipl-vn}

\title{Tập luận văn đội tuyển quốc gia Trung Quốc năm 2018}
\author{Huấn luyện viên: Zhang Ruizhe}
\date{China Computer Federation, tháng 4 năm 2018}

\begin{document}
\maketitle

\origtitle{Tập luận văn ứng viên đội tuyển quốc gia Trung Quốc Olympic Tin học năm 2018}
\sourcepath{IOI China candidate team papers / National training team 2018 collection.pdf}

\renewcommand{\abstractname}{Tóm tắt}
\begin{abstract}
PDF nguồn là tập hợp luận văn đội tuyển quốc gia Trung Quốc năm 2018. Khác với
một số tập trước, lớp văn bản của tệp này trích xuất được khá tốt bằng
\texttt{pdftotext}: phần bìa, mục lục và các trang mở đầu của từng bài đều đọc
được. Bản dịch này chuyển ngữ phần đầu sách và mục lục, đồng thời ghi lại các
chủ đề chính có thể xác nhận từ mục lục và phần mở đầu của các bài. Bản hiện
tại đã bổ sung bản dịch đầy đủ hai bài đầu tiên: bài về hàm sinh trong các bài
toán gieo xúc xắc, và báo cáo ra đề \emph{Số nút của cây hậu tố}.
\end{abstract}

\section{Thông tin đầu sách}

Tập sách có tiêu đề \emph{Tập luận văn ứng viên đội tuyển quốc gia Trung Quốc
Olympic Tin học năm 2018}. Huấn luyện viên ghi trên bìa là Zhang Ruizhe. Đơn vị
xuất bản là China Computer Federation; thời điểm ghi trên bìa là tháng 4 năm
2018.

\section{Mục lục đã dịch}

\begin{center}
\small
\begin{tabular}{r p{0.46\linewidth} p{0.29\linewidth}}
\textbf{Trang} & \textbf{Bài viết} & \textbf{Tác giả / trường} \\
\hline
1 & Bàn sơ lược về ứng dụng hàm sinh trong bài toán gieo xúc xắc & Yang Maolong, Trường Trung học Changjun, Trường Sa \\
11 & Báo cáo ra đề \emph{Số nút của cây hậu tố} và tối ưu hóa một lớp bài toán đoạn & Chen Jianglun, Trường Trung học Changjun, Trường Sa \\
23 & Bàn sơ lược về bài toán hồi quy bảo toàn thứ tự & Gao Ruiquan, Trường Ngoại ngữ Nam Kinh \\
34 & Báo cáo ra đề \emph{Fim 4} & Wu Jinzhao, Trường Trung học số 2 Hàng Châu, Chiết Giang \\
45 & Một số kỹ thuật và công cụ giải bài toán khối liên thông trên cây & Ren Xuandi, Trường Trung học số 1 Thiệu Hưng \\
58 & Báo cáo ra đề \emph{Jellyfish} và khảo cứu mở rộng & Liang Yancheng, Trường Trung học Trấn Hải, Ninh Ba \\
75 & \emph{Leafy Tree} và cây cân bằng có trọng số được hiện thực bằng nó & Wang Siqi, Trường Trung học số 7 Thành Đô \\
87 & Báo cáo ra đề \emph{Tiểu H thích tô màu} & Chen Jiale, Trường Trung học số 2 Hàng Châu, Chiết Giang \\
92 & Một số bài toán tổng hàm số học đặc biệt & Zhu Zhenting, THPT trực thuộc Đại học Sư phạm An Huy \\
113 & Bàn sơ lược về ứng dụng DFT trong thi tin học & Liu Chengao, THPT trực thuộc Đại học Công nghiệp Tây Bắc \\
132 & Báo cáo ra đề \emph{Hàng đợi hoàn hảo} & Lin Xuheng, THPT trực thuộc Đại học Sư phạm Phúc Kiến \\
143 & Bàn sơ lược về một số mở rộng của matroid và ứng dụng & Yang Qianlan, Trường Trung học Hoài Âm, Giang Tô \\
164 & Bàn sơ lược về tính chất của Splay và Treap cùng ứng dụng & Dong Weijun, Trường Trung học số 6 Quảng Châu \\
180 & Báo cáo ra đề \emph{Cây khung phương sai nhỏ nhất} & He Zhongtian, Trường Trung học số 80 Bắc Kinh \\
192 & Khảo cứu các bài toán sinh và đếm liên quan tới đồ thị Euler & Chen Tong, Trường Thực nghiệm trực thuộc Đại học Sư phạm Bắc Kinh
\end{tabular}
\end{center}

\section{Bài 1: Bàn sơ lược về ứng dụng hàm sinh trong bài toán gieo xúc xắc}

\begin{center}
Yang Maolong, Trường Trung học Changjun, Trường Sa
\end{center}

\subsection*{Tóm tắt}

Bài toán gieo xúc xắc là một lớp bài toán thường gặp trong thi thuật toán. Sau
khi nghiên cứu kỹ lớp bài toán này, tác giả nhận thấy chúng đều có thể được
giải bằng hàm sinh. Vì vậy, bài viết này xoay quanh chuỗi bài toán gieo xúc
xắc, trình bày cách dùng hàm sinh để từng bước giải quyết vấn đề, cũng như ưu
thế của cách giải bằng hàm sinh so với phương pháp truyền thống.

\noindent\textbf{Từ khóa:} hàm sinh, bài toán gieo xúc xắc, xác suất và kỳ vọng.

\subsection*{1 Dẫn nhập}

Bài toán gieo xúc xắc là một lớp bài toán thường gặp trong thi thuật toán, và
đã nhiều lần xuất hiện trong các kỳ thi những năm gần đây.

Mặt khác, hàm sinh là một công cụ hữu hiệu để giải lớp bài toán này. So với
phương pháp truyền thống, nó có các ưu điểm như dễ tính toán và khả năng mở
rộng mạnh. Tuy nhiên, hiện nay trong giới OI vẫn có rất ít nghiên cứu về phương
pháp này.

Mục 2 quy ước một số ký hiệu. Mục 3 giới thiệu định nghĩa của hàm sinh xác
suất và một số tính chất. Mục 4 giới thiệu ứng dụng cơ bản của thuật toán này.
Mục 5 và 6 kết hợp với các đề bài để giới thiệu một số ứng dụng phức tạp hơn.

\subsection*{2 Quy ước ký hiệu}

\begin{enumerate}
  \item \(A_i\) biểu thị số thứ \(i\) của dãy \(A\).
  \item \(A[l,r]\) biểu thị dãy gồm các số thứ \(l\) đến thứ \(r\) của dãy
  \(A\).
  \item \(f^{(k)}(x)\) là đạo hàm cấp \(k\) của \(f(x)\).
  \item \([P]\) là ngoặc Iverson: nếu điều kiện trong ngoặc vuông đúng thì giá
  trị bằng \(1\), nếu không thì bằng \(0\).
\end{enumerate}

\subsection*{3 Kiến thức chuẩn bị}

\subsubsection*{3.1 Hàm sinh thông thường}

Hàm sinh thông thường của dãy \(a_0,a_1,a_2,\ldots\) là
\[
  A(x)=\sum_{i=0}^{\infty} a_i x^i .
\]

\subsubsection*{3.2 Hàm sinh xác suất}

Nếu với dãy \(a_0,a_1,a_2,\ldots\) tồn tại một biến ngẫu nhiên rời rạc \(X\) nào
đó thỏa \(\Pr(X=i)=a_i\), thì hàm sinh thông thường của \(a_n\)
\((n\in\mathbb{N})\) được gọi là hàm sinh xác suất của \(X\).

Nói cách khác, nếu \(X\) là biến ngẫu nhiên rời rạc trên tập số nguyên không âm
\(\mathbb{N}\), thì hàm sinh xác suất của \(X\) là
\[
  F(z)=\mathbb{E}(z^X)=\sum_{i=0}^{\infty}\Pr(X=i)z^i .
\]

\subsubsection*{3.2.1 Tính chất của hàm sinh xác suất}

Vì \(X\) là biến ngẫu nhiên rời rạc trên \(\mathbb{N}\), ta nhất định có
\[
  F(1)=\sum_{i=0}^{\infty}\Pr(X=i)=1 .
\]

Lấy đạo hàm \(F(z)\), ta được
\[
  F'(z)=\sum_{i=0}^{\infty} i\Pr(X=i)\cdot z^{i-1}.
\]

Do đó kỳ vọng của \(X\) là
\[
  \mathbb{E}(X)=F'(1)=\sum_{i=0}^{\infty} i\Pr(X=i).
\]

Suy luận thêm có thể nhận được
\[
  \mathbb{E}(X^k)=F^{(k)}(1),\quad k\ne 0 .
\]

Vì vậy phương sai của \(X\) là
\[
  \operatorname{Var}(X)=F''(1)+F'(1)-\bigl(F'(1)\bigr)^2 .
\]

\subsubsection*{3.3 Border}

Với một dãy \(A\) có độ dài \(L\), nếu
\[
  A[1,i]=A[L-i+1,L],
\]
thì gọi \(A[1,i]\) là một \emph{border} của \(A\).

\subsection*{4 Dạng bài cơ bản}

\paragraph{Ví dụ 1. Vương quốc ca hát}
\noindent\emph{Nguồn: CTSC2006.}

Cho một dãy \(A\) có độ dài \(L\). Sau đó, mỗi lần gieo một con xúc xắc công
bằng ghi các số từ \(1\) đến \(m\), rồi đưa số gieo được vào cuối dãy \(B\) ban
đầu rỗng. Nếu trong dãy \(B\) đã xuất hiện dãy \(A\) cho trước, tức \(A\) là một
dãy con liên tiếp của \(B\), thì dừng lại. Hãy tính độ dài kỳ vọng của dãy
\(B\). \(L\le 10^5\).

\subsubsection*{4.1 Phân tích}

Dạng cơ bản của bài toán gieo xúc xắc là: cho một dãy và một con xúc xắc, sau đó
liên tục gieo xúc xắc để sinh ra một dãy ngẫu nhiên, cho đến khi dãy đã cho
xuất hiện trong dãy ngẫu nhiên thì dừng, và cần tính số lần gieo kỳ vọng.

\subsubsection*{4.2 Một phương pháp không dùng hàm sinh}

Đặt \(E_i\) là số lần gieo xúc xắc kỳ vọng từ trạng thái hiện tại thỏa
\(A[1,i]\) là hậu tố của \(B\) cho đến lần đầu tiên thỏa \(A[1,i+1]\) là hậu tố
của \(B\). Khi đó đại lượng cần tìm là
\[
  \sum_{i=0}^{n-1} E_i .
\]

Định nghĩa \(f_{i,j}\) là độ dài border dài nhất không phải chính nó của dãy mới
thu được bằng cách thêm số \(j\) vào sau \(A[1,i]\). Khi đó có phương trình
\[
  E_i=1+\frac{1}{m}\sum_{j=1}^{m}[j\ne A_{i+1}]
      \sum_{k=f_{i,j}}^{i}E_k .
\]

Giải phương trình có thể nhận được
\[
  E_i=m+(m-1)\sum_{k=1}^{i-1}E_k
      -\sum_{j=1}^{m}[j\ne A_{i+1}]
       \sum_{k=1}^{f_{i,j}-1}E_k .
\]

Độ phức tạp thời gian như vậy là \(O(nm)\). Vấn đề chính nằm ở việc tính
\[
  \sum_{j=1}^{m}[j\ne A_{i+1}]
  \sum_{k=1}^{f_{i,j}-1}E_k,
\]
và phần này có thể tối ưu. Dựa trên công thức so sánh \(A[1,i]\) với border dài
nhất không phải chính nó, có thể thấy giữa hai công thức chỉ có \(O(1)\) giá trị
\(f_{i,j}\) khác nhau, nên mỗi lần chỉ cần bảo trì lại \(O(1)\) giá trị này. Khi
đó độ phức tạp thời gian giảm xuống \(O(n)\).

Có thể thấy phương pháp này rất phức tạp và khả năng mở rộng thấp. Phương pháp
hàm sinh được giới thiệu dưới đây ưu việt hơn phương pháp này khá nhiều.

\subsubsection*{4.3 Phương pháp hàm sinh}

Đặt \(a_i\) biểu thị \(A[1,i]\) có phải là một border của \(A\) hay không:
\(a_i=1\) nghĩa là có, còn \(a_i=0\) nghĩa là không.

Đặt \(f_i\) là xác suất để khi kết thúc, độ dài dãy ngẫu nhiên bằng \(i\); hàm
sinh xác suất của nó là \(F(x)\). Đặt dãy phụ trợ \(g_i\) là xác suất để độ dài
dãy ngẫu nhiên đạt đến \(i\) mà vẫn chưa kết thúc; hàm sinh thông thường của nó
là \(G(x)\).

Đại lượng cần tìm là \(F'(1)\). Qua phân tích, ta nhận được các đẳng thức sau:
\begin{align}
  F(x)+G(x) &= 1+G(x)\cdot x, \tag{1}\\
  G(x)\cdot\left(\frac{1}{m}x\right)^L
    &= \sum_{i=1}^{L}a_i\cdot F(x)\cdot
       \left(\frac{1}{m}x\right)^{L-i}. \tag{2}
\end{align}

Ý nghĩa của (1) là: sau khi thêm một số vào một dãy chưa kết thúc, dãy có thể
kết thúc hoặc vẫn chưa kết thúc; \(1\) là trường hợp ban đầu khi dãy ngẫu nhiên
rỗng.

Ý nghĩa của (2) là: sau khi thêm dãy đã cho vào sau một dãy chưa kết thúc, chắc
chắn sẽ kết thúc, nhưng có thể đã kết thúc trước khi thêm đủ \(L\) số; khi đó
dãy đã thêm nhất định phải là một border của dãy đã cho.

Lấy đạo hàm (1) rồi thay \(x=1\), ta được
\begin{align}
  F'(x)+G'(x) &= G'(x)\cdot x+G(x),\\
  F'(1) &= G(1). \tag{3}
\end{align}

Tiếp tục thay \(x=1\) vào (2), ta được
\begin{align}
  G(1)\cdot\left(\frac{1}{m}\right)^L
    &= \sum_{i=1}^{L}a_i\cdot F(1)\cdot
       \left(\frac{1}{m}\right)^{L-i},\\
  G(1)&=\sum_{i=1}^{L}a_i\cdot m^i . \tag{4}
\end{align}

Từ (3)(4), suy ra
\[
  F'(1)=\sum_{i=1}^{L}a_i\cdot m^i .
\]

Dùng thuật toán KMP hoặc hash là có thể tìm \(a_i\) và tính \(F'(1)\) trong độ
phức tạp thời gian \(O(L)\).

\subsubsection*{4.4 Suy nghĩ thêm}

Ta có thể tiến thêm một bước để nhận được phương pháp tính phương sai.

Theo
\[
  \operatorname{Var}(X)=F''(1)+F'(1)-\bigl(F'(1)\bigr)^2,
\]
có thể thấy chỉ cần biết \(F'(1)\) và \(F''(1)\). \(F'(1)\) đã biết, vì vậy cần
tính \(F''(1)\).

Lấy đạo hàm cấp hai của (1) rồi thay \(x=1\), ta được
\begin{align}
  F''(x)+G''(x) &= G'(x)+G''(x)\cdot x+G'(x),\\
  F''(1) &= 2G'(1).
\end{align}

Tiếp tục lấy đạo hàm (2) rồi thay \(x=1\), ta được
\begin{align}
  G(x) &= \sum_{i=1}^{L}a_i\cdot m^i\cdot F(x)\cdot x^{-i},\\
  G'(x) &= \sum_{i=1}^{L}a_i\cdot m^i\cdot
          \bigl(F'(x)\cdot x^{-i}-i\cdot F(x)\cdot x^{-i-1}\bigr),\\
  G'(1) &= \sum_{i=1}^{L}a_i\cdot m^i\cdot \bigl(F'(1)-i\bigr).
\end{align}

Vì vậy
\[
  F''(1)=2\cdot\sum_{i=1}^{L}a_i\cdot m^i\cdot \bigl(F'(1)-i\bigr).
\]

Nếu tiếp tục suy luận và quy nạp, có thể nhận được
\[
  F^{(k)}(1)=k\cdot\sum_{i=1}^{L}a_i\cdot m^i
  \sum_{j=0}^{k-1}(-1)^j\binom{k-1}{j}
  \frac{(i+j-1)!}{(i-1)!}F^{(k-1-j)}(1).
\]

\subsubsection*{4.5 Tiểu kết}

Phương pháp dùng hàm sinh để giải lớp bài toán này thường bắt đầu bằng cách
định nghĩa một hàm sinh xác suất \(F(x)\) và một hàm sinh phụ trợ \(G(x)\). Sau
đó xét việc thêm một số hoặc một dãy cho trước vào trạng thái chưa kết thúc, rồi
dựa vào tình huống thực tế để lập phương trình. Cuối cùng, thông qua thay giá
trị và lấy đạo hàm, ta giải ra \(F'(1)\) cần tìm.

\subsection*{5 Dạng bài phức tạp}

\paragraph{Ví dụ 2. Trò chơi gieo xúc xắc}
Cho \(n\) dãy \(A_i\) có độ dài lần lượt là \(L_i\). Lại cho một con xúc xắc ghi
các số từ \(1\) đến \(m\), trong đó xác suất gieo ra \(i\) là \(P_i\). Sau đó,
mỗi lần gieo xúc xắc một lần và đưa số trên xúc xắc vào cuối dãy \(B\) ban đầu
rỗng. Nếu cả \(n\) dãy đã cho đều là dãy con liên tiếp của \(B\), thì dừng lại.
Hãy tính độ dài kỳ vọng của dãy \(B\). Bảo đảm \(n\le 15\), \(L\le 20000\),
\(m\le 10^5\).

\subsubsection*{5.1 Phân tích}

Về bản chất, bài toán này vẫn là bài toán gieo xúc xắc, nhưng đã biến thành
trường hợp cần \(n\) dãy đều xuất hiện trong dãy ngẫu nhiên.

\subsubsection*{5.2 Lời giải}

Đặt \(T_i\) là độ dài của dãy ngẫu nhiên khi \(A_i\) xuất hiện lần đầu. Khi đó
đại lượng cần tìm là
\[
  \mathbb{E}\!\left(\max_{i\in\{1,2,\ldots,n\}}T_i\right).
\]

Nếu một dãy \(A_i\) là dãy con liên tiếp của một dãy khác \(A_j\), thì nhất định
có \(T_i\le T_j\), nên có thể loại bỏ \(A_i\).

Giá trị lớn nhất không dễ tính, nên dùng bao hàm--loại trừ Min--Max:
\[
  \max_{i=1}^{n}T_i
  =\sum_{S\subseteq\{1,2,\ldots,n\}}(-1)^{|S|+1}
   \min_{i\in S}T_i .
\]

Theo tính tuyến tính của kỳ vọng, bài toán có thể chuyển thành tính
\(\mathbb{E}(\min_{i\in S}T_i)\), tức hễ một dãy xuất hiện thì kết thúc. Không
mất tính tổng quát, giả sử \(S=\{1,2,\ldots,n\}\).

Định nghĩa
\[
  P(A)=\prod_{i\in A}P_i .
\]
Đặt
\[
  a_{i,j,k}=[A_i[1,k]=A_j[L_j-k+1,L_j]] .
\]

Đặt \(f_{i,j}\) là xác suất để dãy xuất hiện đầu tiên là \(A_i\) và độ dài dãy
ngẫu nhiên bằng \(j\); hàm sinh thông thường của nó là \(F_i(x)\). Đặt dãy phụ
trợ \(g_i\) là xác suất để độ dài dãy ngẫu nhiên đạt đến \(i\) mà vẫn chưa kết
thúc; hàm sinh thông thường của nó là \(G(x)\).

Ta có các đẳng thức sau:
\begin{align}
  \sum_{i=1}^{n}F_i(x)+G(x) &= 1+G(x)\cdot x, \tag{5}\\
  G(x)\cdot P(A_i)x^{L_i}
    &= \sum_{j=1}^{n}\sum_{k=1}^{L_i}
       a_{i,j,k}\cdot F_j(x)\cdot P(A_i[k+1,L_i])\cdot x^{L_i-k},
       \quad \forall i\in S . \tag{6}
\end{align}

Đại lượng cần tính là
\[
  \sum_{i=1}^{n}F_i'(1).
\]

Lấy đạo hàm (5) rồi thay \(x=1\), ta được
\[
  \sum_{i=1}^{n}F_i'(1)=G(1).
\]

Thay \(x=1\) vào (6), ta được
\[
  G(1)=\sum_{j=1}^{n}
       \left(\sum_{k=1}^{L_i}a_{i,j,k}\cdot\frac{1}{P(A_i[1,k])}\right)
       \cdot F_j(1),\quad \forall i\in S .
\]

Hiện có \(n+1\) ẩn, nhưng chỉ có \(n\) phương trình. Có thể thấy
\(\sum_{i=1}^{n}F_i(x)\) là hàm sinh xác suất của độ dài dãy ngẫu nhiên, nên
\[
  \sum_{i=1}^{n}F_i(1)=1 .
\]
Như vậy có thể dùng khử Gauss để tìm \(G(1)\). Đồng thời cũng có thể nhận được
\(F_i(1)\), tức xác suất để dãy xuất hiện đầu tiên là dãy thứ \(i\); chẳng hạn
bài \([{\rm SDOI2017}]\) Trò chơi đồng xu yêu cầu tính \(F_i(1)\).

Dùng hash, có thể tìm mảng \(a\) trong độ phức tạp thời gian \(O(n^2L)\), và
với mỗi cặp \((i,j)\) tính
\[
  \sum_{k=1}^{L_i}a_{i,j,k}\cdot\frac{1}{P(A_i[1,k])}.
\]
Với mỗi \(S\), có thể dùng khử Gauss trong \(O(n^3)\) để giải \(G(1)\). Vì vậy
tổng độ phức tạp thời gian là \(O(n^2L+2^n n^3)\).

\subsection*{6 Một số ứng dụng mới}

\subsubsection*{6.1 Khớp mơ hồ}

\paragraph{Ví dụ 3. Dice}
\noindent\emph{Nguồn: 2013 Multi-University Training Contest 5.}

Với một con xúc xắc công bằng \(m\) mặt, hãy tính số lần gieo kỳ vọng để dừng
khi \(n\) kết quả cuối cùng giống nhau, hoặc để dừng khi \(n\) kết quả cuối cùng
đôi một khác nhau. Bảo đảm \(n,m\le 10^6\), và với câu hỏi thứ hai bảo đảm
\(n\le m\).

\subsubsection*{6.1.1 Phương pháp sai phân truyền thống}

\paragraph{Câu hỏi thứ nhất}
Đặt \(f_i\) là số lần gieo xúc xắc kỳ vọng từ trạng thái hiện tại đến khi kết
thúc, trong đó trạng thái hiện tại thỏa \(i\) lần cuối cùng giống nhau. Khi đó
đại lượng cần tìm là \(f_0\).

Có thể lập công thức chuyển:
\[
  f_i=\frac{1}{m}f_{i+1}+\frac{m-1}{m}f_1+1 .
\]

Trước hết sai phân, ta được
\[
  f_{i-1}-f_i=\frac{1}{m}(f_i-f_{i+1}).
\]

Vì \(f_0-f_1=1\), nên \(f_{i-1}-f_i=m^{i-1}\). Đồng thời \(f_n=0\), vì vậy
\[
  f_0=\sum_{i=0}^{n-1}m^i=\frac{m^n-1}{m-1}.
\]

\paragraph{Câu hỏi thứ hai}
Đặt \(g_i\) là số lần gieo xúc xắc kỳ vọng từ trạng thái hiện tại đến khi kết
thúc, trong đó trạng thái hiện tại thỏa \(i\) lần cuối cùng đôi một khác nhau.
Khi đó đại lượng cần tìm là \(g_0\).

Công thức chuyển là
\[
  g_i=\frac{m-i}{m}g_{i+1}+\frac{1}{m}\sum_{j=1}^{i}g_j+1 .
\]

Trước hết sai phân:
\[
  g_{i-1}-g_i=\frac{m-i}{m}(g_i-g_{i+1}).
\]

Tương tự có \(g_0-g_1=1\), \(g_n=0\), nên cũng có thể tìm \(g_0\) trong thời
gian \(O(n)\).

\subsubsection*{6.1.2 Cách làm bằng hàm sinh}

\paragraph{Phân tích}
Tuy có thể xem đơn giản bài này là một bài toán gieo xúc xắc nhiều dãy, số lượng
dãy hợp lệ lại quá lớn. Nhưng tất cả chúng đều thỏa cùng một điều kiện, vì vậy
có thể gom chúng lại với nhau.

\paragraph{Câu hỏi thứ nhất}
Đặt \(f_i\) là xác suất để kết thúc ở lần thứ \(i\); hàm sinh xác suất của nó là
\(F(x)\). Đặt \(g_i\) là xác suất để đã gieo \(i\) lần mà vẫn chưa kết thúc; hàm
sinh thông thường của nó là \(G(x)\).

Khi đó nhận được
\begin{align}
  F(x)+G(x)&=G(x)\cdot x+1, \tag{7}\\
  G(x)\cdot\left(\frac{1}{m}x\right)^{n-1}
    &=\sum_{i=1}^{n}F(x)\cdot\left(\frac{1}{m}x\right)^{n-i}. \tag{8}
\end{align}

Áp dụng cách giải tương tự phần trước, ta có thể giải ra
\[
  F'(1)=\frac{m^n-1}{m-1}.
\]

\paragraph{Câu hỏi thứ hai}
Đặt \(f_i\) là xác suất để kết thúc ở lần thứ \(i\); hàm sinh xác suất của nó là
\(F(x)\). Đặt \(g_i\) là xác suất để đã gieo \(i\) lần mà vẫn chưa kết thúc; hàm
sinh thông thường của nó là \(G(x)\).

Khi đó nhận được
\begin{align}
  F(x)+G(x)&=G(x)\cdot x+1, \tag{9}\\
  G(x)\cdot\left(\frac{1}{m}x\right)^n\cdot\frac{m!}{(m-n)!}
    &=\sum_{i=1}^{n}F(x)\cdot
      \left(\frac{1}{m}x\right)^{n-i}\cdot
      \frac{(m-i)!}{(m-n)!}. \tag{10}
\end{align}

Có thể giải ra
\[
  F'(1)=\sum_{i=1}^{n}m^i\frac{(m-i)!}{m!}.
\]

\subsubsection*{6.1.3 So sánh}

Tuy hai phương pháp đều có thể cho cùng kết quả, cách làm bằng hàm sinh rõ ràng
đơn giản hơn về mặt tính toán so với phương pháp truyền thống, và phương pháp
cũng dễ hiểu hơn. Đồng thời, với các đề bài khác nhau, phần cần sửa đổi là rất
ít.

\subsubsection*{6.2 Trường hợp phức tạp hơn}

\paragraph{Ví dụ 4. Một bài xác suất nhỏ thú vị}
\noindent\emph{Nguồn: \url{http://roosephu.github.io/2017/12/31/condexp/}.}

Một con xúc xắc \(n\) mặt, mỗi mặt được đánh số từ \(1\) đến \(n\). Có một bộ
đếm ban đầu bằng \(0\). Mỗi lần tung xúc xắc, nếu kết quả là số lẻ thì bộ đếm bị
xóa về \(0\); nếu không thì bộ đếm tăng thêm \(1\), và nếu kết quả là \(n\) thì
kết thúc. Hỏi kỳ vọng của bộ đếm khi kết thúc. Bảo đảm \(n\) là số chẵn.

\subsubsection*{6.2.1 Lời giải}

Nhìn qua thì đáp án là \(\frac{n}{2}\), nhưng thật ra không phải.

Tương tự, đặt \(f_i\) là xác suất để kết thúc khi bộ đếm bằng \(i\); hàm sinh
xác suất của nó là \(F(x)\). Nhưng định nghĩa của mảng phụ trợ cần đổi một
chút: đặt \(g_i\) là số lần kỳ vọng để bộ đếm bằng \(i\); hàm sinh thông thường
của nó là \(G(x)\).

Ta có thể nhận được các phương trình
\begin{align}
  F(x)+G(x)&=G(x)\cdot\frac{x}{2}+\frac{G(1)}{2}+1, \tag{11}\\
  F(x)&=G(x)\cdot\frac{x}{n}. \tag{12}
\end{align}

Ý nghĩa của phương trình thứ nhất là: với mỗi trạng thái của \(G\), có xác suất
\(\frac{1}{2}\) để bộ đếm tăng thêm \(1\), và có xác suất \(\frac{1}{2}\) để bộ
đếm bị xóa về \(0\).

Vì \(F(1)=1\), nên \(G(1)=n\). Phương trình (11) có thể biến thành
\[
  F(x)+G(x)=G(x)\cdot\frac{x}{2}+\frac{n}{2}+1. \tag{13}
\]

Thay (12) vào (13), ta được
\begin{align}
  G(x)\cdot\left(1+\frac{x}{n}-\frac{x}{2}\right)&=\frac{n}{2}+1,\\
  G(x)&=\frac{\frac{n}{2}+1}{1+\frac{x}{n}-\frac{x}{2}},\\
  G(x)&=\frac{n^2+2n}{2n-(n-2)x}. \tag{14}
\end{align}

Lại thay (14) ngược vào (12), ta biết
\begin{align}
  F(x)&=\frac{(n+2)x}{2n-(n-2)x},\\
  F'(x)&=\frac{2n\cdot(n+2)}{[2n-(n-2)x]^2},\\
  F'(1)&=\frac{2n}{n+2}.
\end{align}

\subsection*{7 Tổng kết}

Với lớp bài toán gieo xúc xắc, ta có thể xây dựng hàm sinh và sử dụng thông tin
trong đề bài để lập quan hệ giữa các hàm sinh, rồi giải ra giá trị mong muốn.
Dùng phương pháp này có thể tránh phương pháp khử Gauss truyền thống.

Lớp bài toán được nhắc đến trong bài viết này vẫn đáng được nghiên cứu sâu hơn.
Vì vậy tác giả hy vọng bài viết có thể đóng vai trò gợi mở, và hy vọng các độc
giả quan tâm có thể nghiên cứu tiếp, mở rộng cách làm của tác giả, từ đó nhận
được nhiều cách làm và ứng dụng thú vị hơn.

\subsection*{Lời cảm ơn}

\begin{itemize}
  \item Cảm ơn China Computer Federation đã cung cấp nền tảng học tập và giao
  lưu.
  \item Cảm ơn sự hướng dẫn của huấn luyện viên đội tuyển quốc gia Zhang
  Ruizhe.
  \item Cảm ơn thầy Xie Qiufeng đã quan tâm và hướng dẫn.
  \item Cảm ơn bạn Chen Jianglun và đàn anh Guo Xiaoxu đã cung cấp ý tưởng cho
  tôi.
  \item Cảm ơn bạn Chen Jianglun, học trưởng Luo Yuping, đàn anh Guo Xiaoxu và
  bạn Wang Xiuhan đã đọc duyệt bài viết này.
  \item Cảm ơn gia đình đã quan tâm và ủng hộ tôi.
\end{itemize}

\subsection*{Tài liệu tham khảo}

\begin{enumerate}
  \item Wikipedia, \emph{Probability-generating function}.
  \item Ronald L. Graham, Donald E. Knuth, Oren Patashnik. \emph{Concrete
  Mathematics}.
  \item Jin Ce, \emph{Phép toán trên hàm sinh và bài toán đếm tổ hợp}, luận văn
  ứng viên đội tuyển quốc gia Trung Quốc Olympic Tin học năm 2015.
  \item Hu Yuanming, \emph{Phân tích cơ sở và ứng dụng của xác suất trong thi
  tin học}, luận văn ứng viên đội tuyển quốc gia Trung Quốc Olympic Tin học năm
  2013.
  \item Wei Taiyun, \emph{Các cách giải khác nhau và so sánh cho một bài toán
  tung đồng xu}.
\end{enumerate}

\section{Bài 2: Báo cáo ra đề \emph{Số nút của cây hậu tố} và tối ưu hóa một lớp bài toán đoạn}

\begin{center}
Chen Jianglun, Trường Trung học Changjun, Trường Sa
\end{center}

\subsection*{Tóm tắt}

Cấu trúc dữ liệu hậu tố là một hệ thống hiện đang được ứng dụng rộng rãi trong
thi thuật toán. Nó bao gồm các nội dung quen thuộc như mảng hậu tố, automaton
hậu tố, cây hậu tố, v.v. Những cấu trúc dữ liệu hậu tố này có thể xử lý hiệu quả
rất nhiều bài toán.

Tác giả đã nghiên cứu các thuật toán liên quan tới cấu trúc dữ liệu hậu tố xuất
hiện trong những năm gần đây, rồi ra đề \emph{Số nút của cây hậu tố}. Dựa vào
tính chất của cây hậu tố và automaton hậu tố, tác giả nhận được một cách làm
sơ bộ. Trong quá trình trao đổi ở bài tập đội tuyển và trại mùa đông, tác giả
đã thảo luận cách làm này với các bạn khác, rồi cùng họ nghiên cứu ra các thuật
toán hiệu quả và ngắn gọn hơn.

Khi nghiên cứu thuật toán cho bài này, tác giả nhận thấy rằng với lớp bài toán
chuỗi dùng cây động để bảo trì cây hậu tố, có thể dùng hợp nhất theo heuristic
để thay cây động, từ đó làm cho mô hình này được cài đặt đơn giản hơn. Đồng
thời, phương pháp này còn có thể mở rộng sang nhiều bài toán đoạn hơn.

Bài viết này trình bày chi tiết nguồn gốc ý tưởng ra đề và quá trình tối ưu
từng bước, đồng thời đưa ra một phương án tối ưu cho một lớp bài toán đoạn. Tác
giả hy vọng bài viết có thể giúp nhiều người hiểu sâu hơn các tính chất liên
quan của cấu trúc dữ liệu hậu tố, cũng như học được phương pháp tối ưu cho một
lớp bài toán đoạn.

\subsection*{1 Dẫn nhập}

Trong một kỳ ACM-ICPC khu vực năm 2016, đã xuất hiện một mô hình mới mẻ dùng
cây động để bảo trì cây hậu tố. Mô hình này có thể giải hiệu quả một lớp bài
toán như nhiều lần hỏi số xâu con phân biệt về bản chất của một xâu con trong
một xâu, hoặc độ dài xâu con dài nhất xuất hiện ít nhất hai lần trong một xâu
con. Điều đó khiến mọi người chú ý tới hướng này.

Trong kỳ tự kiểm tra đội tuyển quốc gia năm 2017, phiên bản tăng cường của bài
toán này xuất hiện: đếm số xâu con phân biệt về bản chất trong một xâu con, với
điều kiện bắt đầu bằng một ký tự cố định và kết thúc bằng một ký tự cố định.

Đồng thời, những năm gần đây cũng thường xuất hiện một dạng bài: đã biết một
thuật toán giải một bài toán nào đó, hãy tính một số thông tin trong quá trình
chạy thuật toán ấy. Chẳng hạn trong ZJOI2017 Day1 T2, đã biết một thuật toán hỗ
trợ cộng điểm đơn và hỏi tổng đoạn; yêu cầu của đề không phải là đáp án mà
thuật toán đó trả về, mà là xác suất thuật toán tính đúng kết quả theo đúng quá
trình của nó.

Cây hậu tố là một công cụ mạnh để giải bài toán chuỗi. Từ đó tác giả nghĩ rằng:
nếu không yêu cầu dùng cây hậu tố để giải bài toán chuỗi nói chung, mà hỏi số
nút nhận được sau khi dựng cây hậu tố cho một xâu thì sao?

Từ đó tác giả bắt đầu nghiên cứu vấn đề này và ra đề bài này.

Cấu trúc bài viết như sau:
\begin{itemize}
  \item Mục 2 đến 4 là một số định nghĩa cơ bản, mô tả thông tin đề bài và tình
  hình điểm số.
  \item Mục 5 đến 10 giới thiệu 3 thuật toán ăn điểm từng phần và 3 thuật toán
  đạt điểm tối đa.
  \item Mục 11 đề xuất phương án dùng hợp nhất theo heuristic để thay cây động,
  qua đó tối ưu độ phức tạp mã nguồn của lớp bài toán đoạn như dùng cây động để
  bảo trì cây hậu tố, đồng thời trình bày một số ứng dụng của phương án này.
\end{itemize}

\subsection*{2 Mô tả đề bài}

Cho một xâu \(P\) độ dài \(n\). Có \(m\) truy vấn; mỗi truy vấn cho hai tham số
\(l,r\), hỏi số nút của cây hậu tố tạo bởi xâu con \(P[l,r]\).

\subsubsection*{2.1 Phạm vi dữ liệu}

Trong bài này dùng các số để biểu thị xâu. Các số khác nhau đại diện cho các ký
tự khác nhau, số giống nhau đại diện cho cùng một ký tự; giá trị các số nằm
trong đoạn \([0,n]\).

\begin{itemize}
  \item Với \(20\%\) dữ liệu: \(n\le 500,\ m\le 500\).
  \item Với \(40\%\) dữ liệu: \(n\le 3000,\ m\le 3000\).
  \item Với \(50\%\) dữ liệu: \(n\times m\le 3\times 10^8\).
  \item Với \(10\%\) dữ liệu khác: xâu toàn là \(0\).
  \item Với \(20\%\) dữ liệu khác: mỗi ký tự là một số nguyên không âm ngẫu nhiên
  trong \([0,n]\), và \(l,r\) của mỗi truy vấn đều được sinh ngẫu nhiên.
  \item Với \(100\%\) dữ liệu: \(n\le 10^5,\ m\le 3\times 10^5\), mỗi số trong
  xâu đều không vượt quá \(n\).
\end{itemize}

\subsubsection*{2.2 Giới hạn tài nguyên}

Giới hạn thời gian \(3\) giây, giới hạn bộ nhớ \(1\) GB.

\subsubsection*{2.3 Định dạng nhập}

Dòng đầu gồm hai số nguyên dương \(n,m\), biểu thị độ dài xâu và số truy vấn.

Dòng thứ hai gồm \(n\) số nguyên không âm cách nhau bởi dấu cách, biểu thị xâu
đó.

Tiếp theo là \(m\) dòng, mỗi dòng gồm hai số nguyên dương \(l,r\) (\(l\le r\)),
biểu thị truy vấn xâu con \([l,r]\).

\subsubsection*{2.4 Định dạng xuất}

In \(m\) dòng, mỗi dòng một số nguyên dương biểu thị đáp án.

\subsubsection*{2.5 Ví dụ nhập}

\begin{verbatim}
7 1
2 1 3 1 3 1 4
1 7
\end{verbatim}

\subsubsection*{2.6 Ví dụ xuất}

\begin{verbatim}
10
\end{verbatim}

\noindent Khi tính bài này, bỏ qua nút gốc.

\subsubsection*{2.7 Giải thích ví dụ}

Cho số \(1\) biểu thị ký tự \(a\), số \(2\) biểu thị ký tự \(b\), số \(3\) biểu
thị ký tự \(n\), số \(4\) biểu thị ký tự \(s\). Khi đó cây hậu tố của xâu trong
ví dụ chính là cây hậu tố của \texttt{bananas}. Hình trong bản gốc vẽ cây này
với các cạnh lần lượt gắn các nhãn như \texttt{bananas}, \texttt{a},
\texttt{na}, \texttt{s}, \texttt{na}, \texttt{s}, \texttt{nas}, v.v., và tính
đáp án sau khi bỏ qua gốc.

\subsection*{3 Phân bố điểm}

Có \(59\) người tham gia lần huấn luyện này. Trong đội tuyển và nhóm bồi dưỡng
tinh anh, có \(4\) người đạt điểm tối đa, \(1\) người đạt \(60\) điểm, \(1\)
người đạt \(50\) điểm, \(2\) người đạt \(40\) điểm, \(2\) người đạt \(30\) điểm,
và \(1\) người đạt \(20\) điểm.

\subsection*{4 Các khái niệm liên quan trong bài viết}

Với một xâu \(S\), dùng \(S[i]\) hoặc \(S_i\) biểu thị ký tự thứ \(i\) của
\(S\), dùng \(S[l,r]\) biểu thị xâu con của \(S\) nằm trong đoạn \([l,r]\), và
dùng \(\lvert S\rvert\) biểu thị độ dài xâu \(S\).

Định nghĩa ký tự kế tiếp của một xâu con \(S[a,b]\) của \(S\) là \(S[b+1]\).
Đặc biệt, nếu \(S[b]\) là cuối xâu thì ký tự kế tiếp của \(S[a,b]\) không được
định nghĩa.

\subsubsection*{4.1 Hash}

Hash là một công cụ thường dùng để xử lý bài toán chuỗi; nó thực hiện một ánh
xạ nén. Trong bài viết này sẽ dùng hash để ánh xạ mỗi xâu vào một số nguyên,
biến bài toán xét hai xâu có bằng nhau hay không thành bài toán xét hai số
nguyên có bằng nhau hay không.

\subsubsection*{4.2 Cây trie hậu tố}

Cây trie hậu tố của một xâu \(S\) có \(\lvert S\rvert\) ký tự là một cây có
hướng chứa nút gốc. Mỗi nút đại diện cho một xâu; các nút khác nhau đại diện
cho các xâu khác nhau; mỗi nút không phải gốc đại diện cho một xâu con của
\(S\); nút gốc biểu thị xâu rỗng. Trên cây, xâu mà một tổ tiên đại diện là một
tiền tố của xâu mà nút đang xét đại diện.

\subsubsection*{4.3 Cây hậu tố}

Cây hậu tố là một dạng nén của cây trie hậu tố.

Trong cây trie hậu tố của \(S\), nếu xâu mà một nút đại diện là một hậu tố của
\(S\), thì gọi nút đó là nút hậu tố.

Trên cơ sở cây trie hậu tố, sau khi gộp mỗi nút không phải nút hậu tố và có
đúng một con với nút con của nó, ta nhận được cây hậu tố.

Trong cây hậu tố, dùng \(fa_k\) biểu thị cha của nút \(k\) trên cây. Dùng
\(last_k\) biểu thị vị trí xuất hiện cuối cùng của xâu mà nút \(k\) đại diện.

\subsubsection*{4.4 Automaton hậu tố}

Automaton hậu tố của một xâu \(S\) là một automaton hữu hạn có thể chấp nhận mọi
hậu tố của \(S\).

Định nghĩa và phương pháp xây dựng automaton hậu tố đã rất phổ biến trong thi
thuật toán, nên bài viết này không trình bày lại. Độc giả chưa quen có thể xem
tài liệu tham khảo [1].

Cấu trúc cây tạo bởi tất cả các nút của automaton hậu tố được gọi là cây
\emph{parent}.

\subsubsection*{4.5 Cây động}

Cây động là một cấu trúc dữ liệu thường dùng trong thi thuật toán hiện nay, dùng
để bảo trì động phân rã đường của một cây.

Nội dung cơ bản của cây động đã rất phổ biến; độc giả chưa quen có thể xem tài
liệu tham khảo [6].

\subsection*{5 Thuật toán một}

Dựa vào cây trie hậu tố để xây dựng cây hậu tố: trước hết dùng mọi hậu tố của
một xâu để xây dựng một cây trie, sau đó gộp mọi nút không phải nút hậu tố và
có đúng một con với nút con của nó.

Độ phức tạp thời gian của thuật toán này là \(O(n^2m)\), độ phức tạp không gian
là \(O(n^3)\). Lý do là kích thước bảng chữ cái bằng \(n\), còn số nút trong cây
trie của mọi hậu tố là \(O(n^2)\).

Điểm kỳ vọng: \(20\) điểm.

\subsection*{6 Thuật toán hai}

Với dữ liệu \(n,m\le 3000\), có thể độc lập tính đáp án của từng xâu trong mỗi
truy vấn với độ phức tạp tương đối tốt.

\paragraph{Bổ đề 6.1.} Cây \emph{parent} của automaton hậu tố là cây hậu tố của
xâu đảo.

Tính chất này thường được dùng trong bài toán chuỗi và cũng được nhắc tới trong
tài liệu tham khảo [2].

Dựa vào tính chất này, có thể đảo xâu rồi xây dựng automaton hậu tố. Việc xây
dựng automaton hậu tố có thể dùng phương pháp tăng dần tuyến tính.

Thuật toán này có độ phức tạp thời gian \(O(nm\log n)\), độ phức tạp không gian
\(O(n)\). Vì kích thước bảng chữ cái của bài này rất lớn, cần dùng cấu trúc dữ
liệu, chẳng hạn \texttt{map} trong C++, để bảo trì các cạnh đi ra của mỗi điểm;
do đó độ phức tạp thời gian có thêm một nhân tử \(\log\).

Điểm kỳ vọng: \(40\) điểm.

\subsection*{7 Thuật toán ba}

Thực ra việc lưu các cạnh chuyển của automaton hậu tố có thể dùng hash để tối
ưu cấu trúc dữ liệu của thuật toán hai. Khi đó đạt được độ phức tạp thời gian
\(O(nm)\), còn độ phức tạp không gian vẫn là \(O(n)\).

Điểm kỳ vọng: \(50\) điểm.

\subsection*{8 Thuật toán bốn}

\subsubsection*{8.1 Phân tích sơ bộ}

Ký hiệu số nút cây hậu tố của xâu cần xét \(S\) là \(A(S)\).

Các nút của cây hậu tố có thể chia làm hai loại: nút hậu tố và nút không phải
hậu tố. Số nút hậu tố đúng bằng \(\lvert S\rvert\); phần này chỉ liên quan tới
độ dài xâu.

Vì nút lá của cây trie hậu tố nhất định là nút hậu tố, nên một nút không phải
hậu tố nhất định có ít nhất một con. Nếu một nút không phải hậu tố có đúng một
con thì trong cây hậu tố nút này sẽ bị gộp với con của nó. Do đó, những nút cần
được tính vào đáp án phải thỏa có ít nhất hai con.

Một nút trên cây trie hậu tố đại diện cho một xâu con của \(S\). Nếu nút \(k\)
có ít nhất hai con, điều đó có nghĩa là tồn tại ít nhất hai ký tự khác nhau
\(c_1,c_2\), sao cho xâu con \(t\) mà nút \(k\) đại diện thỏa: \(tc_1\) và
\(tc_2\) cũng đều là xâu con của \(S\).

\subsubsection*{8.2 Cách làm phức tạp}

Sau khi phát hiện tính chất trên, tác giả nhận được một cách làm có độ phức tạp
thời gian \(O(n\log^3 n + m\log^2 n)\). Tuy nhiên, cách làm này quá phức tạp,
hiệu suất cũng thấp, rất khó hiểu. Đồng thời phương pháp này không ảnh hưởng
tới việc hiểu các cách làm phía sau, nên bài viết này không trình bày lại. Độc
giả quan tâm có thể tìm phần giới thiệu chi tiết về cách làm này trong tài liệu
tham khảo [5].

Điểm kỳ vọng: \(100\) điểm.

\subsection*{9 Thuật toán năm}

\subsubsection*{9.1 Ứng dụng mô hình cây động bảo trì cây hậu tố}

Điểm khó của bài này nằm ở nhiều nhóm truy vấn, mỗi nhóm hỏi một xâu con
\(P[l,r]\) của xâu nhập \(P\). Xét việc liệt kê đầu trái \(l\) của truy vấn từ
phải sang trái; với mỗi \(r\), ta bảo trì động số nút cây hậu tố của xâu
\(P[l,r]\).

Khi \(l\) dịch sang trái một vị trí, với mỗi xâu \(P[l,r]\), xét các thay đổi
giữa \(P[l,r]\) và \(P[l+1,r]\).

Dựa vào tính chất vừa phát hiện, xét một tiền tố \(P[l,v]\) của \(P[l,r]\), gọi
xâu này là \(t\). Nếu trong \([l,r]\) tồn tại hai ký tự \(c_1,c_2\) sao cho
\(tc_1\) và \(tc_2\) đều là xâu con của \(P[l,r]\), còn xâu \(P[l+1,r]\) không
thỏa điều kiện này, thì xâu \(t\) sẽ tạo thêm đóng góp \(1\) cho đáp án của
\(P[l,r]\). Sau khi đầu trái chuyển từ \(l+1\) sang \(l\), cần lập tức tìm các
\(v\) thỏa yêu cầu này.

Vì \(tc_1\) và \(tc_2\) đều là xâu con của \(P[l,r]\), nên xâu \(t\) chắc chắn
đã xuất hiện trong \(P[l+1,r]\). Do trong \(P[l+1,r]\), \(t\) chưa tạo đóng góp
cho đáp án, nên với mọi vị trí xuất hiện \(P[a,b]\) của \(t\) trong
\(P[l+1,r]\) (\(b-a+1=\lvert t\rvert,\ P[a,b]=t\)), hoặc \(b=r\), hoặc các ký
tự kế tiếp của những xâu này, tức \(P[b+1]\), đều giống nhau.

Hình trong bản gốc minh họa trường hợp trên bằng các lần xuất hiện của \(t\)
trong đoạn \([l,r]\), kèm các ký tự kế tiếp \(c_1,c_2,c_2\).

Mọi tiền tố \(P[l,v]\) (\(l\le v\le r\)) của xâu \(P[l,r]\) tương ứng trên cây
trie hậu tố với một đường đi từ gốc tới một nút nào đó. Giả sử các nút trên
đường này lần lượt là \(x_1,x_2,\ldots,x_{r-l+1}\). Cần thống kê có bao nhiêu
nút \(x_i\) thỏa \(x_i\) tồn tại một con khác \(x_{i+1}\). Nếu tồn tại, vì
\(x_{i+1}\) là con của \(x_i\), nên \(x_i\) đã có ít nhất hai con.

Mỗi lần thêm mọi tiền tố bắt đầu ở \(l\) hiện tại, ta đánh dấu đã thăm cho tất
cả \(x_1,x_2,\ldots,x_{r-l+1}\). Định nghĩa \emph{con thực} của một nút là nút
con được thăm cuối cùng trong tất cả các con của nó. Nếu sau một thao tác, con
thực của một nút \(k\) không thay đổi, điều đó cho thấy ký tự kế tiếp của xâu
\(t\) mà nút \(k\) đại diện giống với ký tự kế tiếp trong lần xuất hiện trước
của xâu này. Trong hình của bản gốc, ký tự kế tiếp của xâu hiện tại
\(t=P[l,l+\lvert t\rvert-1]\) và ký tự kế tiếp của lần xuất hiện trước của
\(t\) đều là \(c_1\).

Chú ý rằng điều kiện để xét một xâu có tạo đóng góp mới cho đáp án hay không
là: hiện tại đã xuất hiện ít nhất hai loại ký tự kế tiếp, còn trước đó chỉ xuất
hiện một loại ký tự kế tiếp. Vì vậy, nếu ký tự kế tiếp của một tiền tố của
\(P[l,r]\) bằng ký tự kế tiếp của lần xuất hiện cuối cùng, thì nó không thể tạo
đóng góp mới cho đáp án.

Nói cách khác, chỉ những nút có con thực thay đổi mới có khả năng tạo đóng góp
mới cho đáp án. Chú ý rằng định nghĩa con thực ở đây giống con thực trong cây
động, và mỗi lần đánh dấu tương đương với thao tác \texttt{access} của cây động.
Theo phân tích độ phức tạp của cây động, số lần chuyển con thực là
\(O(\log n)\) khấu hao, vì vậy điều này liên hệ với mô hình cổ điển dùng cây
động để bảo trì cây hậu tố.

Cụ thể, dùng \(last_t\) biểu thị vị trí xuất hiện trước của xâu \(t\), và dùng
\(diff_t\) biểu thị vị trí xuất hiện trước của xâu \(t\) với ký tự kế tiếp khác.
Hình trong bản gốc minh họa \(last[t]\) và \(diff[t]\) bằng bốn lần xuất hiện
của \(t\) có các ký tự kế tiếp \(c_1,c_2,c_2,c_3\).

Với xâu con \(P[l+1,r]\), khi \(r\ge diff_t+\lvert t\rvert\), \(t\) tạo đóng góp
cho đáp án của \(P[l+1,r]\); khi \(r<diff_t+\lvert t\rvert\), \(t\) chưa tạo
đóng góp cho đáp án.

Với xâu con \(P[l,r]\), khi \(r\ge last_t+\lvert t\rvert\), \(t\) tạo đóng góp
cho đáp án của \(P[l,r]\); khi \(r<last_t+\lvert t\rvert\), \(t\) chưa tạo đóng
góp cho đáp án.

Do đó, khi chuyển từ \(P[l+1,r]\) sang \(P[l,r]\), nếu \(r\) nằm trong đoạn
\([last_t+\lvert t\rvert,\ diff_t+\lvert t\rvert)\), đáp án sẽ mới tăng thêm
\(1\). Đây là một thao tác cộng đoạn.

Như vậy, có thể dùng cây động và cây đoạn để hỗ trợ \texttt{access}, tìm các nút
có con thực thay đổi, và thực hiện thao tác cộng đoạn.

\subsubsection*{9.2 Tiểu kết}

Nội dung trên chính là ứng dụng của một mô hình cổ điển trong bài này.

Nhưng đó chỉ là một phần của bài. Tiếp theo còn phải đối mặt với một vấn đề
khác.

\subsubsection*{9.3 Xử lý đáp án bị tính lặp}

Khi tính số nút cây hậu tố, ta chia toàn bộ nút thành hai loại: nút hậu tố và
nút không phải hậu tố. Phần trước đã thống kê ``nút hậu tố + nút trên cây trie
hậu tố có ít nhất hai con''. Trong đó, các nút vừa là nút hậu tố vừa có hai con
đã bị tính hai lần. Theo nguyên lý bù trừ, chỉ cần trừ số nút đồng thời thỏa
hai điều kiện này là nhận được đáp án cuối cùng.

Khi hỏi xâu \(P[l,r]\), nếu tồn tại một \(v\) sao cho xâu \(P[v,r]\) xuất hiện
trong \(P[l,r]\) với ít nhất hai ký tự kế tiếp khác nhau, thì xâu này sẽ bị tính
lặp.

\subsubsection*{9.4 Tính chất then chốt}

Sau khi phân tích, có thể phát hiện một tính chất then chốt: nếu một xâu \(t\)
có hai ký tự kế tiếp khác nhau trong \(P[l,r]\), thì mọi hậu tố của \(t\) đều có
ít nhất hai ký tự kế tiếp khác nhau.

Dựa vào tính chất này, độ dài các xâu bị tính lặp thỏa tính đơn điệu, nên có
thể dùng chặt nhị phân để tính độ dài lớn nhất của xâu bị tính lặp.

Sau khi chặt nhị phân, bài toán chuyển thành: cho một xâu, hãy xét xâu này có
hai ký tự kế tiếp khác nhau trong một đoạn hay không.

\subsubsection*{9.5 Hợp nhất cây đoạn}

Có hai ký tự kế tiếp khác nhau nghĩa là trên cây trie hậu tố có hai con khác
nhau. Vì vậy cần hỏi một nút có tồn tại hai cây con của hai con chứa nút hậu tố
nằm trong \([l,r]\) hay không.

Vì con thực là nút con được thăm cuối cùng trong các con của một nút, nếu xâu
mà con thực đại diện không xuất hiện trong đoạn \([l,r]\), điều đó cho thấy xâu
mà nút hiện tại đại diện không có bất kỳ ký tự kế tiếp nào trong \([l,r]\), chắc
chắn không thỏa điều kiện. Nếu xâu mà con thực đại diện xuất hiện trong đoạn
\([l,r]\), tiếp theo cần xét trong cây con của các con khác có tồn tại một nút
hậu tố nằm trong \([l,r]\) hay không; vấn đề này có thể giải bằng hợp nhất cây
đoạn.

Dùng hợp nhất cây đoạn để bảo trì tập vị trí bắt đầu của các xâu do nút hậu tố
trong mỗi cây con đại diện, ta có thể truy vấn trong \(O(\log n)\) số nút hậu tố
của một cây con nằm trong đoạn \([l,r]\). Xét ngoài con thực ra còn có cây con
của con nào khác tồn tại nút hậu tố nằm trong \([l,r]\) hay không, tương đương
với truy vấn xem hiệu giữa số nút hậu tố trong cây con của nút hiện tại nằm
trong \([l,r]\) và số nút hậu tố trong cây con của con thực nằm trong \([l,r]\)
có bằng \(0\) hay không.

Trên đây là toàn bộ nội dung của thuật toán năm. Thuật toán dùng cây động + cây
đoạn và chặt nhị phân + hợp nhất cây đoạn; tổng độ phức tạp là
\(O(n\log^2 n + m\log^2 n)\).

Điểm kỳ vọng: \(100\) điểm.

\subsection*{10 Thuật toán sáu}

Bài này thực ra còn có thể tối ưu thêm. Chú ý rằng quá trình xử lý là liệt kê
đầu trái \(l\) của truy vấn, rồi tính đáp án cho mọi truy vấn có đầu trái hiện
tại. Khi \(l\) cố định, với một xâu \(t\), chỉ cần tìm một \(r\) nhỏ nhất sao
cho \(t\) xuất hiện trong \([l,r]\) với ký tự kế tiếp khác nhau. Giá trị \(r\)
này chính là \(diff_{node(t)}+\lvert t\rvert\) đã được bảo trì trước đó, trong
đó \(node(t)\) là nút tương ứng với xâu \(t\) trên cây hậu tố.

Sau khi chặt nhị phân một xâu, có thể dùng bảng hash để truy vấn trong \(O(1)\)
vị trí của một xâu trên cây hậu tố. Như vậy có thể phán đoán trong \(O(1)\) xem
một xâu có xuất hiện trong một đoạn với ký tự kế tiếp khác nhau hay không.

Nhờ đó độ phức tạp thời gian được tối ưu thành \(O(n\log^2 n + m\log n)\).

Điểm kỳ vọng: \(100\) điểm.

\subsection*{11 Tối ưu hóa một lớp bài toán đoạn}

Khi nghiên cứu mô hình cây động bảo trì cây hậu tố, tác giả phát hiện bản chất
của việc chuyển con thực là: với một xâu \(t\), tìm mọi vị trí xuất hiện của nó
và sắp xếp thành một dãy \(pos_1,pos_2,\ldots,pos_v\) theo thứ tự tăng dần của
đầu phải vị trí xuất hiện. Với hai phần tử kề nhau bất kỳ \(pos_i\) và
\(pos_{i+1}\) của dãy này, nếu ký tự kế tiếp của hai vị trí này khác nhau, điều
đó cho thấy các nút lá hậu tố tương ứng với hai xâu này nằm trong cây con của
hai con khác nhau của nút tương ứng với \(t\) trên cây hậu tố. Khi đầu trái truy
vấn được liệt kê từ \(pos_{i+1}\) tới \(pos_i\), con thực trên cây hậu tố sẽ
chuyển một lần.

Một khi đã bảo trì dãy như vậy, ta cũng có thể tính được \(last\) và \(diff\)
mà không cần cây động.

Dùng hợp nhất theo heuristic để tính dãy \(pos\) của mỗi cây con. Đồng thời,
hợp nhất theo heuristic cũng có thể chứng minh độ phức tạp số lần chuyển con
thực trong bài này. Độ phức tạp chính là số lần xuất hiện tình huống
\(pos_i\) và \(pos_{i+1}\) có ký tự kế tiếp khác nhau.

Định nghĩa con nặng là con có số nút hậu tố trong cây con lớn nhất trong tất cả
các con, các con khác là con nhẹ. Nếu trước hết chỉ xét con nặng, ký tự kế tiếp
của \(pos_1\) tới \(pos_v\) đều giống nhau. Sau đó thêm từng nút hậu tố của con
nhẹ vào; mỗi lần thêm một nút vào dãy \(pos\), nhiều nhất tạo ra \(2\) điểm đứt
mới với hai phần tử kề trước và sau. Vì mỗi số nhiều nhất được chèn \(O(\log n)\)
lần với tư cách thuộc con nhẹ, nên tổng độ phức tạp là \(O(n\log n)\) nhân với
độ phức tạp của phép cộng đoạn.

Do ngay từ bước đầu tác giả đã nghĩ tới việc dùng hợp nhất theo heuristic để
thay cây động, nên đã không nhận ra rằng phần sau có thể trực tiếp dùng phương
pháp trong thuật toán sáu để phán đoán. Dĩ nhiên, vì hợp nhất theo heuristic
cũng có thể tính được \(diff\) và \(last\) tương ứng, thuật toán sáu vẫn có thể
được cài đặt bằng hợp nhất theo heuristic.

Phương pháp như vậy không chỉ dùng được cho bài này, mà còn có thể dùng trong
nhiều bài toán khác.

\subsubsection*{11.1 Ví dụ một}

\paragraph{11.1.1 Mô tả đề bài và phạm vi dữ liệu}

Cho một xâu \(P[1,n]\), có \(m\) truy vấn; mỗi truy vấn hỏi số xâu con phân biệt
về bản chất trong một xâu con \([l,r]\).

\[
  n,m\le 10^5 .
\]

\paragraph{11.1.2 Phân tích}

Dựa vào ý tưởng được cung cấp trong bài viết này, liệt kê đầu trái \(l\) của
truy vấn từ phải sang trái, rồi bảo trì đáp án cho mỗi đầu phải.

Vẫn xét lượng thay đổi của đáp án khi đầu trái truy vấn chuyển từ \(l+1\) thành
\(l\). Với một xâu \(P[l,v]\), nếu đầu phải của vị trí xuất hiện trước của xâu
này lớn hơn \(r\), còn đầu phải \(v\) của vị trí xuất hiện hiện tại không vượt
quá \(r\), thì nó tạo đóng góp mới \(1\) cho đáp án. Tương tự, dùng \(last_t\)
biểu thị vị trí xuất hiện cuối cùng của xâu \(t\). Khi đó với mọi
\(v\in[l,n]\), đáp án của các truy vấn có đầu phải nằm trong
\([v,last_{P[l,v]}-1]\) cần tăng thêm \(1\).

Tiếp theo, nếu xâu \(P[l,v]\) và xâu \(P[l,v+1]\) có \(last\) giống nhau, thì có
thể xử lý chung, tương đương với cộng một cấp số cộng có công sai \(1\) lên một
đoạn.

Theo tính chất của cây động, mọi điểm nằm trên cùng một đường thực có thời điểm
được thăm cuối cùng giống nhau; tức là \(last\) của các xâu mà những nút này đại
diện giống nhau. Do đó, các điểm trên cùng một đường thực có thể được xử lý
cùng lúc. Mỗi lần \texttt{access}, số đường thực được thăm là \(O(\log n)\) khấu
hao, nên tổng độ phức tạp thời gian là \(O(n\log^2 n)\), độ phức tạp không gian
\(O(n)\).

Khi \texttt{access}, ta sẽ tìm được mọi đường thực từ một nút tới gốc. Điểm đổi
giữa hai đường thực chính là nơi xảy ra một lần chuyển con thực. Phần trước đã
nói rằng hợp nhất theo heuristic có thể truy vấn mọi sự kiện chuyển con thực,
nên có thể dùng hợp nhất theo heuristic để thay cây động.

\subsubsection*{11.2 Tiểu kết}

Các bài toán chuỗi dùng được tối ưu này thường có các đặc điểm sau:
\begin{enumerate}
  \item Có một xâu mẫu lớn được nhập vào, mỗi truy vấn hỏi thông tin của một xâu
  con của xâu này.
  \item Có thể biết toàn bộ xâu mẫu lớn trước khi trả lời truy vấn, không phải
  bắt buộc trực tuyến thêm ký tự vào cuối xâu mẫu.
  \item Đáp án cần tìm liên quan tới mỗi xâu con phân biệt về bản chất và tiền
  nhiệm, hậu nhiệm của nó trong xâu mẫu.
\end{enumerate}

\subsubsection*{11.3 Mở rộng}

Khi nghiên cứu sâu hơn quá trình này, tác giả phát hiện phương pháp này có thể
mở rộng sang nhiều bài toán đoạn hơn.

\subsubsection*{11.4 Ví dụ hai}

\paragraph{11.4.1 Mô tả đề bài và phạm vi dữ liệu}

Cho một cây \(n\) nút, có \(m\) nhóm truy vấn. Mỗi lần hỏi số nút sau khi xây
dựng cây ảo trên tất cả các nút có số hiệu nằm trong \([l,r]\).

\[
  n,m\le 10^5 .
\]

\paragraph{11.4.2 Phân tích}

Với truy vấn \([l,r]\), xét xem mỗi điểm \(k\) có tạo đóng góp cho đáp án hay
không. \(k\) nằm trong cây ảo nếu \(l\le k\le r\), hoặc nếu trong cây con của ít
nhất hai con của \(k\) đều tồn tại một nút nằm trong đoạn \([l,r]\).

Vẫn dùng ý tưởng tương tự: liệt kê đầu trái \(l\) từ phải sang trái, rồi bảo trì
động đáp án cho mỗi đầu phải \(r\). Mỗi lần thực hiện một thao tác
\texttt{access} với đầu trái hiện tại \(l\); việc chuyển cạnh ảo-thực tương ứng
với việc đáp án của các truy vấn có đầu phải \(r\) trong một đoạn sẽ mới tăng
thêm \(1\).

Cũng có thể dùng hợp nhất theo heuristic: sắp xếp mọi nút trong cây con của
\(k\) theo số hiệu. Khi \([l,r]\) chứa hai nút đến từ cây con của hai con khác
nhau, \(k\) sẽ nằm trong cây ảo. Ở đầu mục này đã chứng minh tổng số đoạn là
\(O(n\log n)\). Vì vậy bài này có thể được giải trong thời gian
\(O(n\log^2 n)+O(m\log n)\).

\subsection*{12 Tổng kết}

Bài này là một bài toán lấy số nút cây hậu tố làm mô hình, khảo sát năng lực
phân tích tính chất cây hậu tố của thí sinh và năng lực vận dụng mô hình cổ
điển. Các cấu trúc liên quan gồm automaton hậu tố, cây hậu tố, cây đoạn, cây
động và hợp nhất theo heuristic, đều thuộc phạm vi kiến thức của thí sinh NOI.

Tuy nhiên, bài này chú trọng kiểm tra khả năng phân tích tính chất của cấu trúc
dữ liệu hậu tố. Chẳng hạn so với cây hậu tố, mỗi nút không bị nén trên cây trie
hậu tố tương ứng với một xâu con trong xâu gốc; xâu con đó hoặc là một hậu tố,
hoặc có hai ký tự kế tiếp khác nhau.

Bài này gồm hai phần, mỗi phần đều cần suy nghĩ sâu mới nhận được phương án
giải hiệu quả; đối với thí sinh, độ khó là rất lớn.

Ban đầu bài này được tạo ra khi nghiên cứu tính chất của automaton hậu tố. Lúc
đầu tác giả đặt trực tiếp vấn đề lên automaton hậu tố và nhận được thuật toán
bốn rất phức tạp. Đưa bài này vào bài tập đội tuyển là để thử tìm cách làm đơn
giản hơn.

Sau đó, bạn Zhu Zhenting từ THPT trực thuộc Đại học Sư phạm An Huy phát hiện
một tính chất quan trọng, tức tính chất được nhắc tới trong bài viết này: độ dài
các xâu bị tính lặp thỏa tính đơn điệu. Qua chặt nhị phân, cách làm được đơn
giản hóa rất nhiều. Trong buổi trao đổi của trại mùa đông năm nay, bạn ấy cùng
bạn Wang Xiuhan từ Trường Trung học số 7 Thành Đô đã chia sẻ ý tưởng này. Đó
chính là thuật toán năm trong bài viết này.

Sau buổi trao đổi, tác giả và bạn Zhang Yubo từ Trường Trung học số 80 Bắc Kinh
đã thảo luận sâu hơn, phát hiện chỉ cần dùng \(last\), \(diff\) và bảng hash là
có thể phán đoán trong \(O(1)\) xem một xâu có từng xuất hiện trong một đoạn với
ký tự kế tiếp khác nhau hay không. Đó chính là thuật toán sáu trong bài viết
này.

Trong quá trình nghiên cứu các thuật toán này, tác giả phát hiện bản chất của
cây động bảo trì cây hậu tố là tìm mọi vị trí xuất hiện của một xâu, rồi chia
chúng theo ký tự kế tiếp: các xâu kề nhau và có ký tự kế tiếp giống nhau được
chia vào cùng một đoạn; giữa các đoạn khác nhau sẽ phát hiện thay đổi của
\(last\) và \(diff\). Vì vậy có thể dùng hợp nhất theo heuristic để thay cây
động, làm đơn giản độ phức tạp mã nguồn. Đồng thời tác giả phát hiện phương
pháp này có thể mở rộng sang một số bài toán đoạn khác.

Hy vọng thông qua bài này, mọi người có thể hiểu sâu hơn về tính chất của
automaton hậu tố và cây hậu tố, đồng thời thành thạo cách giải tinh tế cho một
lớp bài toán đoạn.

\subsection*{13 Lời cảm ơn}

\begin{itemize}
  \item Cảm ơn China Computer Federation đã cung cấp nền tảng học tập và giao
  lưu.
  \item Cảm ơn sự hướng dẫn của huấn luyện viên đội tuyển quốc gia Zhang
  Ruizhe.
  \item Cảm ơn thầy Xie Qiufeng của Trường Trung học Changjun đã quan tâm và
  hướng dẫn.
  \item Cảm ơn Wang Xiuhan, Zhu Zhenting và Zhang Yubo đã giúp đỡ về ý tưởng
  tối ưu của bài này.
  \item Cảm ơn Luo Yuping, Guo Xiaoxu, Yang Maolong, Liu Chengao, Zhang
  Qingchuan và Wu Hang đã đọc duyệt bài viết này.
  \item Cảm ơn gia đình đã quan tâm và ủng hộ tôi.
\end{itemize}

\subsection*{Tài liệu tham khảo}

\begin{enumerate}
  \item Chen Lijie, \emph{Automaton hậu tố}, trao đổi trại mùa đông năm 2012.
  \item Zhang Tianyang, \emph{Automaton hậu tố và ứng dụng}, tập luận văn ứng
  viên đội tuyển quốc gia Trung Quốc năm 2015.
  \item Hong Huadun, \emph{Báo cáo ra đề Tái cấu trúc hệ gene}, tập luận văn ứng
  viên đội tuyển quốc gia Trung Quốc năm 2017.
  \item \emph{String}, ACM/ICPC Asia Regional Qingdao Online 2016.
  \item Chen Jianglun, \emph{Báo cáo lời giải đề tự chọn trong bài tập đội tuyển
  quốc gia IOI2018}.
  \item Huang Zhi'ao, \emph{Bàn sơ lược về các vấn đề liên quan tới cây động và
  một số mở rộng đơn giản}, tập luận văn ứng viên đội tuyển quốc gia Trung Quốc
  năm 2014.
\end{enumerate}

\section{Chủ đề chính}

Từ mục lục và phần mở đầu của các bài, tập năm 2018 bao phủ các nhóm chủ đề
sau:
\begin{itemize}
  \item hàm sinh, xác suất, kỳ vọng và các bài toán gieo xúc xắc;
  \item cấu trúc dữ liệu chuỗi, cây hậu tố, suffix automaton, KMP và tính chu
  kỳ của chuỗi;
  \item hồi quy bảo toàn thứ tự, chia đôi toàn cục và tối ưu trên quan hệ thứ
  tự bộ phận;
  \item quy hoạch động và cấu trúc dữ liệu cho khối liên thông trên cây;
  \item tích chập, DFT, FFT, FWT, FMT và các bài toán đa thức;
  \item cây nhị phân, \emph{Leafy Tree}, Splay, Treap, tìm kiếm theo ngón tay và
  cây cân bằng có trọng số;
  \item tổng hàm số học, hàm nhân tính, tổng trên số nguyên tố và các biến thể
  số học đặc biệt;
  \item matroid, tối ưu tổ hợp, cây khung phương sai nhỏ nhất và đồ thị Euler;
  \item các báo cáo ra đề cho \emph{Số nút của cây hậu tố}, \emph{Fim 4},
  \emph{Jellyfish}, \emph{Tiểu H thích tô màu}, \emph{Hàng đợi hoàn hảo} và
  \emph{Cây khung phương sai nhỏ nhất}.
\end{itemize}

\section*{Ghi chú về nguồn}

Tệp PDF có 210 trang. \texttt{pdftotext} trích xuất được Unicode đúng cho phần
đầu sách, mục lục và các trang thân bài đã kiểm tra, nên không cần render mục
lục bằng ảnh để đọc lại. Bản dịch đã bổ sung toàn văn bài đầu tiên của Yang
Maolong, tương ứng trang in 1--10 và trang PDF 3--12, cùng bài thứ hai của Chen
Jianglun, tương ứng trang in 11--22 và trang PDF 13--24. Một số công thức dài,
ví dụ mẫu và hình minh họa được đối chiếu thêm bằng ảnh render từ PDF vì bố cục
công thức và chú thích hình có chỗ bị méo trong văn bản trích xuất.

\end{document}
